% User-defined commands go here.

% GENERAL CUSTOM COMMANDS %%%%%%%%%%%%%%%%%%%%%%%%%%%%%%%%%%%%%%%%%%%%%%%%%%%%%%
% These should work no matter what packages are used (unless the identifiers
% conflict).
%
% Our research tends to quote from interviews a lot, and it can be annoying to
% try and change how all of the quotes are formatted, so set a single format
% here.
\newcommand{\say}[1]{``#1''}
\newcommand{\paraphrase}[1]{[\emph{#1}]}
%
% Similarly, if identifying themes in text, it's helpful to format those in a
% consistent way.
\newcommand{\theme}[1]{\textbf{#1}}
%
% Improve formatting of paragraphs.
%\renewcommand{\paragraph}[1]{\medskip\noindent\textbf{#1}.}
%%%%%%%%%%%%%%%%%%%%%%%%%%%%%%%%%%%%%%%%%%%%%%%%%%%%%%%%%%%%%%%%%%%%%%%%%%%%%%%%

% PACKAGE-SPECIFIC CUSTOMIZATIONS %%%%%%%%%%%%%%%%%%%%%%%%%%%%%%%%%%%%%%%%%%%%%%
% Below are new commands that only work if a specific package is used. For each
% section, begin with a comment line that starts with "Required: ..." to clearly
% mark that it's a package-specific customization.
%
% Required: cite (possibly with noadjust option)
% Fix ugly citation style: [1],[2] --> [1,2].
%\renewcommand{\citepunct}{, }
%\renewcommand{\citedash}{--}
%
% Required: enumitem
% Environment to define commonly-used list types.
% - The default enumerated list has quite a bit of spacing, so make it more
%   compact.
\newenvironment{compactenum}[1][\arabic*.]
{\begin{enumerate}[label={#1}, noitemsep]}
{\end{enumerate}}
% - Same with the itemized list.
\newenvironment{compactitem}%
{\begin{itemize}[noitemsep]}
{\end{itemize}}
% - Customize a list that appears in the flow of a paragraph like (1) one, (2)
%   two, (3) three, (4) four.
\newenvironment{inlineenum}[1][(\arabic*)]
{\begin{enumerate*}[label={#1}, itemjoin={{, }}]}
{\end{enumerate*}}
% - Same as inlineenum, but adds an "and" before the last element to flow more
%   nicely with a sentence, like (1) one, (2) two, and (3) three.
\newenvironment{inparaenum}[1][(\arabic*)]
{\begin{enumerate*}[label={#1}, itemjoin={{, }}, itemjoin*={{, and }}]}
{\end{enumerate*}}
%
% Required: hyperref
% The \autoref command ensures that the proper term is used for the relevant
% reference. Customize the relevant terms for each reference type here.
%\renewcommand*{\sectionautorefname}{Sec.}
%\renewcommand*{\figureautorefname}{Fig.}
%\renewcommand*{\tableautorefname}{Tab.}
% In some templates, subsections are referred to as "Subsection x.y", while
% they may be referred to as "Section x.y" instead. The following line ensures
% that sections and subsections are always referred to with the same identifier.
%\let\subsectionautorefname\sectionautorefname%
%
% Required: msc
% Message sequence charts by default display "msc" in the corner of the chart.
% Turn that off here.
%\setmsckeyword{}
%
% Required: xcolor
% It can be helpful to include formatting for draft text or different authors'
% comments, each with a different text color. Define those here, and also ensure
% that they all get stripped out upon submission.
\ifsubmit%
  \newcommand{\draft}[1]{}
\else
  \newcommand{\draft}[1]{\textcolor{blue}{[Draft text: #1]}}
\fi
%
% Required: xspace
% For some papers, we define new notation that only works in math mode. Ensure
% that these can be easily mentioned in text, with the proper spacing, using the
% following wrapper command.
%\newcommand{\emxs}[1]{\ensuremath{#1}\xspace}
%
%%%%%%%%%%%%%%%%%%%%%%%%%%%%%%%%%%%%%%%%%%%%%%%%%%%%%%%%%%%%%%%%%%%%%%%%%%%%%%%%
